% =========================================================================
% Capítulo 3: Cinética de sistemas de partículas
% Tema 01: Sistemas de partículas
% Problema 10: Péndulo con masa deslizante y resorte
% =========================================================================
\subsubsection{Péndulo con masa deslizante y resorte}

Una partícula de masa $m_1$ está unida al extremo de una barra rígida ideal sin masa de longitud $L$, la cual está articulada en el punto $O$ y puede rotar libremente en el plano vertical. Una segunda masa $m_2$ (cuenta o collarín) puede deslizarse libremente a lo largo de la barra lisa. Esta masa $m_2$ está conectada al pivote $O$ mediante un resorte lineal de constante elástica $k$ y longitud natural $L_0$. Sabiendo que la gravedad actúa hacia abajo, realice lo siguiente:

\begin{enumerate}
    \item Defina un conjunto completo e independiente de coordenadas generalizadas para describir la configuración del sistema.
    \item Analice la cinemática del sistema determinando los vectores de posición y velocidad de ambas masas.
    \item Derive las ecuaciones diferenciales del movimiento aplicando principios de momentum (cantidad de movimiento y momento angular).
\end{enumerate}

\begin{figure}[!ht]
    \centering
    \begin{tikzpicture}[>=Latex, scale=1.4]
        % --- Parámetros ---
        \def\ang{-50}      % Ángulo de la barra
        \def\Lrod{3.5}     % Longitud L
        \def\Rbead{1.8}    % Posición de m2
        
        % Coordenada Origen
        \coordinate (O) at (0,0);

        % --- 1. Techo y Pivote ---
        \fill[pattern=north east lines] (-1.0, 0) rectangle (1.5, 0.3);
        \draw[thick] (-1.0, 0) -- (1.5, 0);
        
        % Pivote O
        \filldraw[thick] (O) circle (3pt);
        \node[below left, yshift=3pt] at (O) {$O$};

        % --- 2. Barra y Resorte ---
        % Barra L
        \draw[line width=5pt] (O) -- (\ang:\Lrod);
        \draw[line width=4pt, gray!30, line cap=round] (O) -- (\ang:\Lrod);
        \draw[thick] (O) -- (\ang:\Lrod);
        \filldraw[thick] (O) circle (0.5pt);

        % Resorte (Desde O hasta m2)
        % Usamos una decoración coil sobre la barra
        \draw[decoration={coil, aspect=0.4, segment length=5pt, amplitude=5pt, pre length=1pt, post length=0pt}, decorate, black] 
            (O) -- (\ang:\Rbead);
        \node[above right, xshift=4pt] at (\ang:{\Rbead/2}) {$k$};

        % --- 3. Masas ---
        % Masa m2 (Cuenta deslizante)
        \coordinate (M2) at (\ang:\Rbead);
        \shadedraw[ball color=gray!30] (M2) circle (0.35) node[] {$m_2$};

        % Masa m1 (Extremo)
        \coordinate (M1) at (\ang:\Lrod);
        \shadedraw[ball color=gray!30] (M1) circle (0.25) node[] {$m_1$};

        % --- 4. Cotas y Gravedad ---
        % Cota L (Longitud total)
        \draw[|<->|] ($(O) + (\ang-90:0.5)$) -- ($(M1) + (\ang-90:0.5)$) node[midway, below left] {$L$};

        % Vector Gravedad
        \draw[->, thick] (2.0, -1.0) -- (2.0, -2.0) node[midway, right] {$g$};

    \end{tikzpicture}
    \caption{Sistema de péndulo compuesto con una masa deslizante sujeta a un resorte.}
    \label{fig:pendulo_masa_deslizante}
\end{figure}